\subsection{诱导测度的性质}

命题 3. --- 设 X 是 T 的一个局部紧子空间，\(\lambda\) 是 X 上的一个复测度。下列性质等价：

\begin{enumerate}
\def\labelenumi{\alph{enumi})}
\item
  典范单射 \(i:\mathrm{X}\to \mathrm{T}\) 是 \(\lambda\) -逆紧的；
\item
  对 T 的每个紧子集 K，\(\mathbf{K} \cap \mathbf{X}\) 本质 \(\lambda\) -可积；
\item
  每个点 \(t \in \mathbf{T}\) 都有一个邻域 \(\mathbf{V}\) ，使得 \(\mathbf{V} \cap \mathbf{X}\) 本质 \(\lambda\) -可积；
\item
  存在 T 上的一个测度 \(\theta\) ，使得 \(\theta_{\mathrm{X}} = \lambda\) 。
\end{enumerate}

若这些等价条件满足，则沿用 d) 中的记号，有

\[
\left(i (\lambda)\right) _ {\mathrm{X}} = \lambda \quad \text{且} \quad i (\lambda) = i (\theta_ {\mathrm{X}}) = \varphi_ {\mathrm{X}} \cdot \theta .\tag{7}
\]

由于单射 \(i\) 是连续的，性质 a)、b)、c) 的等价性由 §6 的命题 1 及其后的注记应用于正测度 \(|\lambda|\) 而得。若 \(\lambda\) 在 X 上由 T 上的一个测度 \(\theta\) 诱导，则 \(|\lambda| = |\theta|_{\mathrm{X}}\) （公式 (6)），于是对 T 的每个紧子集 K 有 \(|\lambda| (\mathrm{K} \cap \mathrm{X}) = |\theta| (\mathrm{K} \cap \mathrm{X}) \leqslant |\theta| (\mathrm{K}) < +\infty\) （命题 1），故 d) 蕴涵 b)。设 a) 满足，让我们证明 \(\left(i(\lambda)\right)_{\mathrm{X}} = \lambda\) ，由此将得 d)。设 \(g\) 是 \(\mathcal{H}(\mathrm{X}; \mathbf{C})\) 的一个元素；用 \(g'\) 表示 \(g\) 到 T 上的零延拓，则依诱导测度的定义，再由 §6 第 4 号的命题 7，有

\[
\int g d (i (\lambda)) _ {X} = \int g ^ {\prime} d (i (\lambda)) = \int (g ^ {\prime} \circ i) d \lambda = \int g d \lambda .
\]

这就完成了四个性质等价性的证明。若 \(\lambda = \theta_{\mathrm{X}}\) 且 \(g\in \mathcal{K}(\mathrm{T};\mathbf{C})\) ，则

\[
\int g d \bigl (i (\theta_ {\mathrm{X}}) \bigr) = \int (g \circ i) d (\theta_ {\mathrm{X}}) = \int g \varphi_ {\mathrm{X}} d \theta ,
\]

因为 \(g\varphi_{\mathbf{X}}\) 是 \(g \circ i\) 到 \(\mathbf{T}\) 上的零延拓。这就证明了 (7) 的第二个公式。

推论 1. --- 若 X 是闭的，则 X 上的每个复测度 \(\lambda\) 都由 T 上的一个测度所诱导。

因为若 K 是 T 中的紧集，则 \(K \cap X\) 是紧的，从而是 \(\lambda\) -可积的。

推论 2. --- 设 \(\theta\) 是 T 上的一个复测度，\(\pi\) 是 T 到局部紧空间 Y 中的一个 \(\theta\) -逆紧映射，\(\pi_{X}\) 是它在 X 上的限制。则 \(\pi_{X}\) 是 \(\theta_{X}\) -逆紧的，且 \(\pi_{X}(\theta_{X}) = \pi(\varphi_{X} \cdot \theta)\) 。

因为 \(\pi_{X} = \pi \circ i\) ，其中 i 是典范单射 \(X \to T\) 。当 \(\theta\) 为正测度时，推论可由命题 3 及像测度的传递性（ \(\S6\) , No.~3, 命题 4）得出。非正复测度的情形随后由线性得出。

命题 4. --- 设 X 和 Y 是 T 的两个局部紧子空间，满足 \(Y \subset X\) 。若 \(\theta\) 是 T 上的一个复测度，则测度 \((\theta_{\mathrm{X}})_{\mathrm{Y}}\) （由 \(\theta_{X}\) 在 Y 上诱导）等于 \(\theta_{Y}\) （“诱导测度的传递性”）。

只须注意：若 g 是 \(\mathcal{K}(\mathrm{Y};\mathbf{C})\) 的一个元素，则 g 到 T 上的零延拓可以由把 g 到 X 上的零延拓再零延拓到 T 而得到；或者，利用概说中的等同，即 \(\varphi_{\mathrm{Y}}\cdot\theta=\varphi_{\mathrm{Y}}(\varphi_{\mathrm{X}}\cdot\theta)\) （§5, No.~4, 命题 8）。

命题 5. --- 设 \((\lambda_{\alpha})_{\alpha\in A}\) 是 T 上的正测度组成的递增有向族，具有上确界 \(\lambda\) ，并设 X 是 T 的一个局部紧子空间。则诱导测度族 \(\lambda_{\alpha}|X\) 在 \(\mathcal{M}(X)\) 中有上界，且

\[
\sup _ {\alpha \in \mathrm{A}} (\lambda_ {\alpha} | \mathrm{X}) = \lambda | \mathrm{X}.\tag{8}
\]

考虑到概说中的等同，本命题是 §5, No.~4 命题 5 的特殊情形。

推论. --- 设 \((\mu_{i})_{i\in\mathbb{I}}\) 是 T 上的正测度组成的可和族，其和为 \(\mu\) 。则诱导测度族 \(\mu_{i}|X\) 是可和的，且

\[
\sum_ {i \in I} (\mu_ {i} | X) = \mu | X.\tag{9}
\]

命题 6. --- 设 \(\Lambda: t \mapsto \lambda_t\) 是一个 \(\mu\) -适宜映射，把 T 映入 \(\mathcal{M}_+(\mathrm{X})\) ，其中 X 是在无穷远处可数的局部紧空间，并设 Y 是 X 的一个局部紧子空间。置 \(\int \lambda_t d\mu(t) = \nu\) 。则映射 \(t \mapsto \lambda_t|\mathrm{Y}\) （T 到 \(\mathcal{M}_+(\mathrm{Y})\) 中的映射）是 \(\mu\) -适宜的，且

\[
\int (\lambda_ {t} | \mathrm{Y}) d \mu (t) = \nu | \mathrm{Y}.\tag{10}
\]

考虑到概说中的等同，本命题是 §5, No.~4 命题 7 的特殊情形。

\section{测度的乘积}

\subsection{乘积测度作为测度积分的解释}

在本节中，T 和 \(T'\) 表示两个局部紧空间，\(\mu\) 是 T 上的一个正测度，\(\mu'\) 是 \(T'\) 上的一个正测度，而 \(\nu = \mu \otimes \mu'\) 是 \(X = T \times T'\) 上的乘积测度（Ch. III, §4, No.~1）。

对每个 \(t \in T\) ，映射 \(t' \mapsto (t, t')\) （\(T'\) 到 X 中）是连续且逆紧的。用 \(\lambda_t'\) 表示 \(\mu'\) 在此映射下的像；\(\lambda_t'\) 是 X 上的一个正测度，且若 \(f \in \mathcal{K}(X)\) ，则用 \(f_t\) 表示偏映射 \(t' \mapsto f(t, t')\) ，有

\[
\int f d \lambda_ {t} ^ {\prime} = \int f _ {t} d \mu^ {\prime},\tag{1}
\]

这也可用关系式 \(\lambda_t' = \varepsilon_t \otimes \mu'\) 来表达。

此外，映射 \(t \mapsto \lambda_t'(f)\) 是连续的且具有紧支集（Ch. III, §4, No.~1, 引理 2），因此映射 \(t \mapsto \lambda_t'\) （T 到 \(\mathcal{M}(\mathbf{X})\) 中的映射）是淡连续的（从而更是淡 \(\mu\) -可测的）；于是测度族 \(t \mapsto \lambda_t'\) 是 \(\mu\) -适宜的（§3, No.~1, 命题 2a)）。\(f\) 关于测度 \(\int \lambda_t' d\mu(t)\) 的积分按定义为

\[
\int \langle f, \lambda_ {t} ^ {\prime} \rangle d \mu (t) = \int d \mu (t) \int f _ {t} (t ^ {\prime}) d \mu^ {\prime} (t ^ {\prime}) = \int f (t, t ^ {\prime}) d \nu (t, t ^ {\prime})
\]

（Ch. III, §4, No.~1, 定理 2）；于是 \(\nu = \int \lambda_t' d\mu(t)\) 。

类似地，对每个元素 \(t^{\prime} \in T^{\prime}\) ，用 \(\lambda_{t^{\prime}}\) 表示 \(\mu\) 在 T 到 X 中的映射 \(t \mapsto (t, t^{\prime})\) 下的像。映射 \(t^{\prime} \mapsto \lambda_{t^{\prime}}\) 是 \(\mu^{\prime}\) -适宜的且淡连续的，且 \(\nu = \int \lambda_{t^{\prime}} d\mu^{\prime}(t^{\prime})\) 。我们将需要下列引理：

引理 1. --- 对每个数值函数 \(f \geqslant 0\) （定义在 \(X\) 上），

\[
\int^ {*} f _ {t} d \mu^ {\prime} = \int^ {*} f d \lambda_ {t} ^ {\prime}.\tag{2}
\]

由于 \(t' \mapsto (t, t')\) 是连续且逆紧的映射，这由 §4, No.~2 的命题 2 得出。

引理 2. --- 设 \(f\) 是 \(X\) 到一个拓扑空间中的映射。为使 \(f\) 是 \(\lambda_t'\) -可测的，必须且只须 \(f_t\) 是 \(\mu'\) -可测的。这是 §6, No.~2 命题 3 的推论。

引理 3. --- 设 f 是定义在 X 上、取值于 \(\overline{R}\) 或一个 Banach 空间中的函数。为使 f 是 \(\lambda_{t}^{\prime}\) -可积的，必须且只须 \(f_{t}\) 是 \(\mu^{\prime}\) -可积的，此时

\[
\int \mathbf {f} _ {t} d \mu^ {\prime} = \int \mathbf {f} d \lambda_ {t} ^ {\prime}.\tag{3}
\]

这由 §4, No.~4 的定理 2 得出，其中用到 \(t' \mapsto (t, t')\) 连续且逆紧这一事实。

注记. --- 引理 1、2、3 可以不用 §§4 和 6 的结果而由一个直接论证非常简单地证明。例如，若 \(f \in \mathcal{K}(\mathrm{T} \times \mathrm{T}')\) ，则关系式 (2) 按定义是明显的。若 \(f\) 在 \(\mathrm{X} = \mathrm{T} \times \mathrm{T}'\) 上下半连续，则只须注意 \(t' \mapsto f_t(t')\) 是诸函数 \(t' \mapsto g_t(t') = g(t, t')\) 的上包络，其中 \(g\) 取遍 \(\mathcal{K}(\mathrm{X})\) 中满足 \(0 \leqslant g \leqslant f\) 的函数所成之集。最后，对任意的 \(f\) ，注意：若 \(h \geqslant f\) 在 \(\mathrm{X}\) 上下半连续，则 \(t' \mapsto h(t, t')\) 在 \(\mathrm{T}'\) 上下半连续；反之，若 \(t' \mapsto u(t')\) 在 \(\mathrm{T}'\) 上下半连续且满足 \(u(t') \geqslant f(t, t')\) （对所有 \(t' \in \mathrm{T}'\) ），则函数 \(h\) ——其中 \(h(t, t') = u(t')\) ，\(h(t_1, t') = +\infty\) （当 \(t_1 \neq t'\) ）——在 \(\mathrm{X}\) 上下半连续且满足 \(h \geqslant f\) 。一旦证明了引理 1，由此即可推出集 \((\mathrm{T} - \{t\}) \times \mathrm{T}'\) 是 \(\lambda_t'\) -可忽略的，于是证明引理 2 和 3 就很容易了。

关系式 (3) 允许把它的两端都记作 \(\int\mathbf{f}(t,t^{\prime})d\mu^{\prime}(t^{\prime})\) 而不致引起混淆。对测度 \(\lambda_{t^{\prime}}=\mu\otimes\varepsilon_{t^{\prime}}\) 显然有类似的结果。

代替记号

\[
\int^ {*} f (t, t ^ {\prime}) d \nu (t, t ^ {\prime}), \quad \int^ {\bullet} f (t, t ^ {\prime}) d \nu (t, t ^ {\prime}), \quad \int \mathbf {f} (t, t ^ {\prime}) d \nu (t, t ^ {\prime}),
\]

我们将采用记号

\[
\iint^ {*} f (t, t ^ {\prime}) d \mu (t) d \mu^ {\prime} (t ^ {\prime}), \iint^ {\bullet} f (t, t ^ {\prime}) d \mu (t) d \mu^ {\prime} (t ^ {\prime}), \iint \mathbf {f} (t, t ^ {\prime}) d \mu (t) d \mu^ {\prime} (t ^ {\prime}),
\]

这与 Ch. III, §4, No.~1 中采用的记号一致。

把测度 \(\nu\) 解释为一个积分 \(\int \lambda_t' d\mu(t)\) ，将使我们可以把 §3 的结果翻译成乘积测度的语言。另一方面，测度 \(\lambda_t'\) 由 \(\{t\} \times T' = \overline{\text{pr}}_1(t)\) 承载，故这个积分定义了 \(\nu\) 的一个切片分解，它是关于投影 \(\text{pr}_1\) （即 \(T \times T'\) 到 \(T\) 上的投影）的（§6, No.~6）。在给出由此所得结果的清单之前，先给出一个有用的性质：

命题 1. --- 设 \((\mu_{\alpha})_{\alpha \in \mathrm{A}}\) （相应地，\((\mu_{\beta}^{\prime})_{\beta \in \mathrm{B}})\) ）是 T 上（相应地，\(\mathbf{T}'\) 上）正测度组成的可和族，其和记作 \(\mu\) （相应地，\(\mu'\) ）。则族 \((\mu_{\alpha} \otimes \mu_{\beta}^{\prime})_{(\alpha, \beta) \in \mathrm{A} \times \mathrm{B}}\) 在 \(\mathbf{T} \times \mathbf{T}'\) 上是可和的，且

\[
\mu \times \mu^ {\prime} = \sum_ {(\alpha , \beta) \in \mathrm{A} \times \mathrm{B}} \mu_ {\alpha} \otimes \mu_ {\beta} ^ {\prime}.\tag{4}
\]

当 A 和 B 有限时，这些性质是明显的。由此推出，若 \(A'\) （相应地，\(B'\) ）是 A（相应地，B）的一个有限子集，则

\[
\sum_ {(\alpha , \beta) \in A ^ {\prime} \times B ^ {\prime}} \mu_ {\alpha} \otimes \mu_ {\beta} ^ {\prime} \leqslant \mu \otimes \mu^ {\prime}.
\]

于是族 \((\mu_{\alpha} \otimes \mu'_{\beta})\) 是可和的。为证明 (4) 的两端相等，只须证明第二端满足乘积测度的特征性质（Ch. III, §4, No.~1, 定理 1），这由下面的计算得出。

设 \(f\) 是 \(\mathcal{H}_{+}(\mathrm{T})\) 的一个元素，\(f'\) 是 \(\mathcal{H}_{+}(\mathrm{T}')\) 的一个元素；回忆 \(f \otimes f'\) 表示函数 \((t, t') \mapsto f(t)f'(t')\) （在 \(\mathrm{T} \times \mathrm{T}'\) 上），它属于 \(\mathcal{H}_{+}(\mathrm{T} \times \mathrm{T}')\) （A, II, §7, No.~7）。于是，由乘积测度的定义，

\[
\begin{array}{r l} \sum_ {(\alpha , \beta) \in \mathrm{A} \times \mathrm{B}} \langle \mu_ {\alpha} \otimes \mu_ {\beta} ^ {\prime}, f \otimes f ^ {\prime} \rangle & = \sum_ {(\alpha , \beta) \in \mathrm{A} \times \mathrm{B}} \big (\langle \mu_ {\alpha}, f \rangle \langle \mu_ {\beta} ^ {\prime}, f ^ {\prime} \rangle \big) \\ & = \Big (\sum_ {\alpha \in \mathrm{A}} \langle \mu_ {\alpha}, f \rangle \Big) \Big (\sum_ {\beta \in \mathrm{B}} \langle \mu_ {\beta} ^ {\prime}, f ^ {\prime} \rangle \Big) \\ & = \langle \mu , f \rangle \langle \mu^ {\prime}, f ^ {\prime} \rangle \\ & = \langle \mu \otimes \mu^ {\prime}, f \otimes f ^ {\prime} \rangle . \end{array}
\]

\subsection{关于两测度乘积可测的函数}

命题 2. --- 设 f 是定义在 T×T′ 上、取值于一个拓扑空间 G 中的 ν-可测函数，并设 M 是由使映射 t′ ↦ f(t, t′) 不是 μ′-可测的 t ∈ T 组成之集。

\begin{enumerate}
\def\labelenumi{\alph{enumi})}
\item
  若 \(f\) 在一个 \(\nu\) -适度子集（\(\mathbf{T} \times \mathbf{T}'\) 的子集）的余集上为常数，则 \(\mathbf{M}\) 是 \(\mu\) -可忽略的。
\item
  若 \(\mu'\) 是适度的，则 M 是 \(\mu\) -可忽略的。
\end{enumerate}

断言 a) 由 §3, No.~2 的命题 4b) 以及 No.~1 的诸注记得出。为处理 b)，注意 \(\mu'\) 是有界测度序列 \(\mu_n'\) 之和（§2, No.~3, 命题 4）；\(f\) 关于 \(\mu \otimes \mu_n' \leqslant \nu\) 可测，且集 M 是诸集 \(M_{n}\) 之并，这些集与诸测度 \(\mu_{n}^{\prime}\) 相伴（§2, No.~2, 命题 2）。于是化归到 \(\mu^{\prime}\) 有界的情形，而这由 §3, No.~2 的命题 4c) 得出。

这一陈述立即推广到复测度（Ch. III, §4, No.~2, 命题 3）。

推论. --- 设 A 是一个 \(\nu\) -可测子集（\(T \times T'\) 的子集），并设 M 是由这样的 \(t \in T\) 组成之集：A 在 t 处的截口 \(\mathrm{A}(t)\) 不是 \(\mu'\) -可测。

\begin{enumerate}
\def\labelenumi{\alph{enumi})}
\item
  若 A 是 \(\nu\) -适度的，则 M 是 \(\mu\) -可忽略的。
\item
  若 A 在 \(T'\) 上的投影是 \(\mu'\) -适度的，则 M 是局部 \(\mu\) -可忽略的。
\end{enumerate}

断言 a) 由命题 2 立即得出。为建立 b)，用 B 表示一个集，它是 T′ 中 \(\mu'\) -可积开集序列之并，且包含 A 在 T′ 上的投影，并用 \(\mu_1'\) 表示适度测度 \(\varphi_{\mathrm{B}} \cdot \mu'\) ；由于 A 关于 \(\mu \otimes \mu_1' \leqslant \mu \otimes \mu'\) 可测，命题 2 蕴涵 A(t) 是 \(\mu_1'\) -可测的（除去组成一个局部 \(\mu\) -可忽略集的 t 外）。但因 A(t) ⊂ B，说 A(t) 是 \(\mu_1'\) -可测的就等价于说 A(t) 是 \(\mu'\) -可测的（§5, No.~3, 命题 4 的推论）。

命题 3. --- 设 f 是 T 到一个拓扑空间 F 中的映射。若 f 是 \(\mu\) -可测的，则映射 \((t, t') \mapsto f(t)\) 是 \(\nu\) -可测的。反之，若 \(\mu' \neq 0\) 且此映射是 \(\nu\) -可测的，则函数 f 是 \(\mu\) -可测的。

第一个断言由 §6, No.~6 命题 10 的推论 1 得出。设 \(\mu' \neq 0\) ，用 \(\mu_1'\) 表示一个被 \(\mu'\) 优超的非零具紧支集测度，用 \(\nu_1\) 表示测度 \(\mu \otimes \mu_1'\) ，并置 \(a = \| \mu_1'\|\) 。于是投影 \(\mathrm{pr}_1\) （\(\mathrm{T} \times \mathrm{T}'\) 到 \(\mathrm{T}\) 上的投影）是 \(\nu_1\) -逆紧的，且像测度 \(\mathrm{pr}_1(\nu_1)\) 等于 \(a\mu\) （§6, No.~6, 命题 10）。若 \((t, t') \mapsto f(t)\) 是 \(\nu\) -可测的，则它也是 \(\nu_1\) -可测的，因此 \(f\) 关于测度 \(a\mu\) 可测（§6, No.~2, 命题 3），再由 \(a \neq 0\) 即得结论。

前述陈述立即推广到复测度（Ch. III, §4, No.~2, 命题 3），下列推论亦然。

推论 1. --- 设 F、F′ 和 G 是三个拓扑空间，u 是 F × F′ 到 G 中的一个连续映射。设 f（相应地，f′）是定义在 T（相应地，T′）上、取值于 F（相应地，F′）中且关于 μ（相应地，μ′）可测的函数。则函数 \((t, t') \mapsto u(f(t), f'(t'))\) 关于 μ ⊗ μ′ 可测。

由于映射 \((t,t^{\prime})\mapsto f(t)\) 、\((t,t^{\prime})\mapsto f^{\prime}(t^{\prime})\) 依命题 3 是 \(\nu\) -可测的，这由 Ch. IV, §5, No.~3 的定理 1 得出。

推论 2. --- 若 A ⊂ T 与 A′ ⊂ T′（分别关于 μ 与 μ′）可测，则 A × A′ 关于 μ ⊗ μ′ 可测。

这由推论 1 立即得出。

推论 3. --- 考虑定义在 T 上的正数值（相应地，复值）函数 f 与 \(f'\) （定义在 \(T'\) 上）。若这两个函数分别关于 \(\mu\) 与 \(\mu'\) 可测，则函数

\[
f \otimes f ^ {\prime}: (t, t ^ {\prime}) \mapsto f (t) f ^ {\prime} (t ^ {\prime})
\]

关于 \(\mu \otimes \mu'\) 可测。

复函数或有限实函数的情形是推论 1 的直接推论。为处理正数值函数的情形，对每个整数 \(n \geqslant 0\) 置 \(f_n = \inf(f, n)\) 、\(f_n' = \inf(f', n)\) ，则有（按惯例 \(0 \cdot (+\infty) = 0\) ）\(f \otimes f' = \sup_n(f_n \otimes f_n')\) ，由此即得结论。

命题 4. --- 设 A 是 T 的一个子集。若 A 是局部 \(\mu\) -可忽略的，则 \(A \times T'\) 是局部 \(\nu\) -可忽略的。反之，若 \(A \times T'\) 是局部 \(\nu\) -可忽略的且 \(\mu' \neq 0\) ，则 A 是局部 \(\mu\) -可忽略的。

第一个断言由 §6, No.~6 命题 10 的推论 1 得出。为建立第二个断言，仍用命题 3 证明中的记号；\(A \times T' = \overline{\text{pr}}_{1}(A)\) 对测度 \(\nu_{1}\) 是局部可忽略的，因此 A 对 \(a\mu\) 是局部可忽略的（§6, No.~2, 命题 2 的推论），再由 \(a \neq 0\) 即得结论。

前述陈述立即推广到两个复测度的乘积（Ch. III, §4, No.~2, 命题 3），下列推论亦然。

推论. --- 若测度 \(\mu\) （相应地，\(\mu'\) ）集中于 M（相应地，M′），则 \(\mu \otimes \mu'\) 集中于 M × M′。

因为 \((\mathrm{T} \times \mathrm{T}') - (\mathrm{M} \times \mathrm{M}')\) 是 \((\mathrm{T} - \mathrm{M}) \times \mathrm{T}'\) 与 \(\mathrm{T} \times (\mathrm{T}' - \mathrm{M}')\) 之并，而依命题 4 这两个集对 \(\mu \times \mu'\) 都是局部可忽略的。

\subsection{正函数的积分}

回忆我们已约定把乘积 \(0 \cdot (+\infty)\) 定义为等于 0。这一约定特别有下列推论：若 f 是数值函数 \(\geqslant 0\) （定义在赋有正测度 \(\lambda\) 的局部紧空间上），则 \(\lambda^{*}(af) = a \cdot \lambda^{*}(f)\) 对每个满足 \(0 \leqslant a \leqslant +\infty\) 的常数 a 成立。当 a = 0 时这是明显的；若 \(a = +\infty\) ，则 \(\lambda^{*}(af) = a \cdot \lambda^{*}(f) = 0\) 或 \(\lambda^{*}(af) = a \cdot \lambda^{*}(f) = +\infty\) （视 f 是否为 \(\lambda\) -可忽略而定）；最后，若 \(0 < a < +\infty\) ，则已知 \(\lambda^{*}(af) = a \cdot \lambda^{*}(f)\) 。

命题 5. --- 设 f 是数值函数 \(\geqslant 0\) （在 \(T \times T'\) 上下半连续）。则函数

\[
t \mapsto \int^ {*} f (t, t ^ {\prime}) d \mu^ {\prime} (t ^ {\prime})
\]

在 T 上是下半连续的，且

\[
\iint^ {*} f (t, t ^ {\prime}) d \mu (t) d \mu^ {\prime} (t ^ {\prime}) = \int^ {*} d \mu (t) \int^ {*} f (t, t ^ {\prime}) d \mu^ {\prime} (t ^ {\prime}).\tag{3}
\]

这是 §3, No.~1 命题 2 的推论，其中用到 No.~1 的引理 1。

推论 1. --- 设 \(f\) （相应地，\(f'\) ）是下半连续函数 \(\geqslant 0\) （定义在 \(T\) （相应地，\(T'\) ）上）；则函数 \(f \otimes f': (t, t') \mapsto f(t)f'(t')\) 在 \(T \times T'\) 上下半连续，且

\[
\iint^ {*} f (t) f ^ {\prime} \left(t ^ {\prime}\right) d \mu (t) d \mu^ {\prime} \left(t ^ {\prime}\right) = \left(\int^ {*} f (t) d \mu (t)\right) \left(\int^ {*} f ^ {\prime} \left(t ^ {\prime}\right) d \mu^ {\prime} \left(t ^ {\prime}\right)\right).\tag{6}
\]

用 \(G\) （相应地，\(G'\) ）表示函数 \(g \in \mathcal{H}_{+}(T)\) （相应地，\(g' \in \mathcal{H}_{+}(T')\) ）——其中 \(g \leqslant f\) （相应地，\(g' \leqslant f'\) ）——所成之集；则

\[
f \otimes f ^ {\prime} = \sup _ {g \in \mathrm{G}, g ^ {\prime} \in \mathrm{G} ^ {\prime}} g \otimes g ^ {\prime}.
\]

由于诸函数 \(g \otimes g'\) 属于 \(\mathcal{H}_{+}(\mathrm{T} \times \mathrm{T}')\) ，\(f \otimes f'\) 确实是下半连续的，且 (6) 由命题 5 立即得出（或直接由在前一公式中取极限得出）。

推论 2. --- 设 A 是 T 的一个 \(\mu\) -适度子集，A′ 是 T′ 的一个 \(\mu'\) -适度子集；则 A × A′ 在 T × T′ 中是 \(\nu\) -适度的。

鉴于适度集的定义（§1, No.~2, 命题 5），只须证明：若 B 是 T 中的可积开集，B′ 是 T′ 中的可积开集，则开集 B × B′ 是可积的。这由推论 1 立即得出。

推论 3. --- 设 A 是 T 的一个 \(\mu\) -可忽略子集，B′ 是 T′ 的一个 \(\mu'\) -适度子集；则 A × B′ 是 \(\nu\) -可忽略的。

因为 \(A \times B'\) 是局部 \(\nu\) -可忽略的（命题 4）且是 \(\nu\) -适度的（推论 2），故是 \(\nu\) -可忽略的（ \(\S1\) , No.~2, 命题 7 的推论 1）。

推论 2 与推论 3 可以推广到两个复测度的乘积，只需把该陈述应用于它们的绝对值（Ch. III, §4, No.~2, 命题 3）。

命题 6. --- 设 \(f\) 是数值函数 \(\geqslant 0\) （定义在 \(\mathbf{T} \times \mathbf{T}'\) 上）。则

\[
\iint^ {*} f (t, t ^ {\prime}) d \mu (t) d \mu^ {\prime} (t ^ {\prime}) \geqslant \int^ {*} d \mu (t) \int^ {*} f (t, t ^ {\prime}) d \mu^ {\prime} (t ^ {\prime}).
\]

这由 §3, No.~2 的命题 3 并注意到 (2) 得出。

命题 7. --- 设 \(f\) 是 \(\nu\) -可测的正数值函数（定义在 \(\mathbf{T} \times \mathbf{T}'\) 上）。

\begin{enumerate}
\def\labelenumi{\alph{enumi})}
\tightlist
\item
  若 \(f\) 是 \(\nu\) -适度的，则函数 \(t \mapsto \int^{*} f(t, t') d\mu'(t')\) 、\(t' \mapsto \int^{*} f(t, t') d\mu(t)\) 分别关于 \(\mu\) 和 \(\mu'\) 可测且适度，且
\end{enumerate}

\[
\begin{array}{r} \iint^ {*} f (t, t ^ {\prime}) d \mu (t) d \mu^ {\prime} (t ^ {\prime}) = \int^ {*} d \mu (t) \int^ {*} f (t, t ^ {\prime}) d \mu^ {\prime} (t ^ {\prime}) \\ = \int^ {*} d \mu^ {\prime} (t ^ {\prime}) \int^ {*} f (t, t ^ {\prime}) d \mu (t). \end{array}\tag{7}
\]

\begin{enumerate}
\def\labelenumi{\alph{enumi})}
\setcounter{enumi}{1}
\tightlist
\item
  若测度 \(\mu'\) 是适度的，则函数 \(t \mapsto \int^{\bullet} f(t, t') d\mu'(t')\) 是 \(\mu\) -可测的，且
\end{enumerate}

\[
\iint^ {\bullet} f (t, t ^ {\prime}) d \mu (t) d \mu^ {\prime} (t ^ {\prime}) = \int^ {\bullet} d \mu (t) \int^ {\bullet} f (t, t ^ {\prime}) d \mu^ {\prime} (t ^ {\prime}).\tag{8}
\]

断言 a) 以及 \(\mu'\) 有界时的断言 b) 都是 §3, No.~2 命题 5 的推论。为处理 \(\mu'\) 适度的情形，把 \(\mu'\) 表示成有界测度序列之和 \(\sum_{n\in\mathbf{N}}\mu_n'\) （§2, No.~3, 命题 4）。于是函数 \(t\mapsto\int^{\bullet}f(t,t')d\mu_n'(t')\) 是 \(\mu\) -可测的，且

\[
\iint^ {\bullet} f (t, t ^ {\prime}) d \mu (t) d \mu_ {n} ^ {\prime} (t ^ {\prime}) = \int^ {\bullet} d \mu (t) \int^ {\bullet} f (t, t ^ {\prime}) d \mu_ {n} ^ {\prime} (t ^ {\prime}).
\]

但 \(\mu \otimes \mu' = \sum_{n \in \mathbf{N}} (\mu \otimes \mu_n')\) （命题 1）；于是对 \(n\) 求和即得断言 b)（ \(\S 2\) , No.~2, 命题 1）。

推论 1. --- 设 H 是 T × T′ 的一个子集，并设 A 是由使 H 在 t 处的截口 H(t) 不是 μ′-可忽略的 t ∈ T 组成之集。

\begin{enumerate}
\def\labelenumi{\alph{enumi})}
\item
  若 H 是 \(\nu\) -可忽略的，则 A 是 \(\mu\) -可忽略的。
\item
  若 H 是局部 \(\nu\) -可忽略的且 \(\mu'\) 是适度的，则 A 是局部 \(\mu\) -可忽略的。
\end{enumerate}

性质 a) 由命题 7（或命题 6）立即得出。在 b) 的假设下，说 \(H(t)\) 局部 \(\mu'\) -可忽略与说它 \(\mu'\) -可忽略是一回事，因为 \(\mu'\) 是适度的（ \(\S1\) , No.~2, 命题 7）。于是性质 b) 由公式 (8) 得出。

这一推论通过取绝对值立即推广到两个复测度的乘积。对下列推论同样成立：

推论 2. --- 若集 \(A \subset T \times T'\) 是 \(\nu\) -可积的，则 A 在 t 处的截口 \(\mathrm{A}(t)\) 是 \(\mu'\) -可积的（对几乎每个 \(t \in T\) ），函数 \(t \mapsto \mu'(\mathrm{A}(t))\) 是 \(\mu\) -可积的，且

\[
\nu (\mathrm{A}) = \int \mu^ {\prime} (\mathrm{A} (t)) d \mu (t).\tag{9}
\]

命题 8. --- 对每对分别定义在 T 和 T′ 上的数值函数 \(f \geqslant 0\) 、\(f' \geqslant 0\) ，

\[
\iint^ {\bullet} f (t) f ^ {\prime} \left(t ^ {\prime}\right) d \mu (t) d \mu^ {\prime} \left(t ^ {\prime}\right) = \left(\int^ {\bullet} f (t) d \mu (t)\right) \left(\int^ {\bullet} f ^ {\prime} \left(t ^ {\prime}\right) d \mu^ {\prime} \left(t ^ {\prime}\right)\right).\tag{10}
\]

我们先处理 \(\mu\) 和 \(\mu'\) 是具紧支集测度的情形；此时 \(\mu \otimes \mu'\) 亦然，且所有符号 \(\int^{\bullet}\) 、\(\iint^{\bullet}\) 都可换成上积分。由命题 6，

\[
\begin{array}{r l} \iint^ {*} f (t) f ^ {\prime} (t ^ {\prime}) d \mu (t) d \mu^ {\prime} (t ^ {\prime}) & \geqslant \int^ {*} d \mu (t) \int^ {*} f (t) f ^ {\prime} (t ^ {\prime}) d \mu^ {\prime} (t ^ {\prime}) \\ & = \left(\int^ {*} f (t) d \mu (t)\right) \left(\int^ {*} f ^ {\prime} (t ^ {\prime}) d \mu^ {\prime} (t ^ {\prime})\right). \end{array}
\]

为建立反向不等式，选取一个函数 \(h \geqslant f\) （相应地，\(h' \geqslant f'\) ），它是一个下半连续函数序列 \((h_n)\) （相应地，\((h'_n)\) ）的下包络，使得

\[
\int^ {*} h (t) d \mu (t) = \int^ {*} f (t) d \mu (t)
\]

\((\text{resp. } \int^{*} h'(t') d\mu'(t') = \int^{*} f'(t') d\mu'(t'))\) ；这样的函数的存在性由上积分的定义（Ch. IV, §1, No.~3, 定义 3）和 Lebesgue 定理立即得出。把命题 7 应用于可测函数 \(h \otimes h'\) ，有

\[
\begin{array}{l} \iint^ {*} f (t) f ^ {\prime} (t ^ {\prime}) d \mu (t) d \mu^ {\prime} (t ^ {\prime}) \leqslant \iint^ {*} h (t) h ^ {\prime} (t ^ {\prime}) d \mu (t) d \mu^ {\prime} (t ^ {\prime}) \\ = \left(\int^ {*} h (t) d \mu (t)\right) \left(\int^ {*} h ^ {\prime} (t ^ {\prime}) d \mu^ {\prime} (t ^ {\prime})\right) \\ = \left(\int^ {*} f (t) d \mu (t)\right) \left(\int^ {*} f ^ {\prime} (t ^ {\prime}) d \mu^ {\prime} (t ^ {\prime})\right), \end{array}
\]

这就是所要的不等式。于是当 \(\mu\) 和 \(\mu'\) 是具紧支集测度时命题得证。为处理一般情形，只须把 \(\mu\) （相应地，\(\mu'\) ）表示成具紧支集测度的族 \((\mu_{\alpha})_{\alpha\in\mathbb{A}}\) （相应地，\((\mu'_{\beta})_{\beta\in\mathbb{B}}\) ）之和（ \(\S2\) , No.~3, 命题 4），对每个测度 \(\mu_{\alpha}\otimes\mu'_{\beta}\) 写出公式 (10)，再对 \((\alpha,\beta)\) 求和，其中用到命题 1（ \(\S2\) , No.~2, 命题 1）。

推论 1. --- 沿用命题 8 中的记号，

\[
\iint^ {*} f (t) f ^ {\prime} \left(t ^ {\prime}\right) d \mu (t) d \mu^ {\prime} \left(t ^ {\prime}\right) = \left(\int^ {*} f (t) d \mu (t)\right) \left(\int^ {*} f ^ {\prime} \left(t ^ {\prime}\right) d \mu^ {\prime} \left(t ^ {\prime}\right)\right)\tag{11}
\]

但可能除去第二端的一个因子等于 0 而另一个等于 +∞的情形。

当第二端的两个因子都有限时，函数 \(f\) 和 \(f'\) 都是适度的（§1, No.~2, 命题 7），因此函数 \(f \otimes f'\) 是适度的（命题 5 的推论 2）；于是上述等式化为公式 (10)（§1, No.~2, 命题 7）。当第二端的一个因子取值 \(+\infty\) 而另一个不为零时，第二端取值 \(+\infty\) ，上述等式由命题 6 得出。

推论 2. --- 设 f 和 \(f'\) 是两个取值于 C 或 \(\overline{R}\) 中的函数，分别定义在 T 和 \(T'\) 上，且分别关于测度 \(\mu\) 和 \(\mu'\) 本质可积（相应地，可积）。则函数 \(f \otimes f'\) 关于测度 \(\mu \otimes \mu'\) 本质可积（相应地，可积），且

\[
\iint f (t) f ^ {\prime} \left(t ^ {\prime}\right) d \mu (t) d \mu^ {\prime} \left(t ^ {\prime}\right) = \left(\int f (t) d \mu (t)\right) \left(\int f ^ {\prime} \left(t ^ {\prime}\right) d \mu^ {\prime} \left(t ^ {\prime}\right)\right).\tag{12}
\]

当 \(f\) 和 \(f'\) 为正时，\(f \otimes f'\) 依命题 3 的推论 3 可测，且该陈述由公式 (10)（相应地，(11)）以及本质可积性判别法（§1, No.~3, 命题 9）（相应地，Ch. IV, §5, No.~6 的可积性判别法，定理 5）得出。一般情形随后立即得出。

推论 2 立即推广到两个复测度的乘积。

\subsection{取值于 Banach 空间中的函数的积分}

定理 1（Lebesgue--Fubini）. --- 设 f 是定义在 T × T′ 上、取值于一个 Banach 空间 F 或 \(\overline{R}\) 中的函数；设 N 是由这样的 \(t \in T\) 组成之集：函数 \(t' \mapsto \mathbf{f}(t, t')\) 不是 \(\mu'\) -可积。

\begin{enumerate}
\def\labelenumi{\alph{enumi})}
\tightlist
\item
  设 \(\mathbf{f}\) 是 \(\nu\) -可积的；则 \(\mathrm{N}\) 是 \(\mu\) -可忽略的，函数 \(t \mapsto \int \mathbf{f}(t, t') d\mu'(t')\) （对 \(t \notin \mathrm{N}\) 有定义）是 \(\mu\) -可积的，且
\end{enumerate}

\[
\iint \mathbf {f} (t, t ^ {\prime}) d \mu (t) d \mu^ {\prime} (t ^ {\prime}) = \int d \mu (t) \int \mathbf {f} (t, t ^ {\prime}) d \mu^ {\prime} (t ^ {\prime}).\tag{13}
\]

\begin{enumerate}
\def\labelenumi{\alph{enumi})}
\setcounter{enumi}{1}
\tightlist
\item
  设 \(\mathbf{f}\) 是本质 \(\nu\) -可积的，且测度 \(\mu'\) 是适度的；则 \(\mathbf{N}\) 是局部 \(\mu\) -可忽略的，函数 \(t \mapsto \int \mathbf{f}(t, t') d\mu'(t')\) （对 \(t \notin \mathbf{N}\) 有定义）是本质 \(\mu\) -可积的，且 (13) 成立。
\end{enumerate}

断言 a) 由 §3, No.~3 的定理 1 立即得出。为建立 b)，用 g 表示一个局部几乎处处等于 f 的 \(\nu\) -可积函数，并用 H 表示由这样的 \((t, t')\) 组成之集：\(\mathbf{f}(t, t') \neq \mathbf{g}(t, t')\) 。依命题 7 的推论 1，截口 \(\mathrm{H}(t)\) 是 \(\mu'\) -可忽略的，除去 \(t \in T\) （组成一个局部 \(\mu\) -可忽略集者）外。于是把 a) 应用于 g 即得关于 f 的结论。

概说. --- 设 f 是定义在 \(T \times T'\) 上、取值于 \(\overline{R}\) 或一个 Banach 空间中的函数，且是 \(\nu\) -可测和 \(\nu\) -适度的。为使三个积分

\[
\iint \mathbf {f} (t, t ^ {\prime}) d \mu (t) d \mu^ {\prime} (t ^ {\prime}), \int d \mu (t) \int \mathbf {f} (t, t ^ {\prime}) d \mu^ {\prime} (t ^ {\prime}), \int d \mu^ {\prime} (t ^ {\prime}) \int \mathbf {f} (t, t ^ {\prime}) d \mu (t)
\]

都存在并且相等，必须且只须两个数

\[
\int^ {*} d \mu (t) \int^ {*} | {\bf f} (t, t ^ {\prime}) | d \mu^ {\prime} (t ^ {\prime}), \quad \int^ {*} d \mu^ {\prime} (t ^ {\prime}) \int^ {*} | {\bf f} (t, t ^ {\prime}) | d \mu (t)
\]

之一为有限。

这是定理 1、命题 7 和可积性判别法（Ch. IV, §5, No.~6, 定理 5）的直接推论。

注记. --- 1) 当测度 \(\mu'\) 不是适度的时候，可能发生这样的情形：\(\mathbf{f}\) 本质 \(\nu\) -可积，而函数 \(t' \mapsto \mathbf{f}(t, t')\) 不是本质 \(\mu'\) -可积的（对任何 \(t \in \mathrm{T}\) ）（§3, 习题 4）。

\begin{enumerate}
\def\labelenumi{\arabic{enumi})}
\setcounter{enumi}{1}
\tightlist
\item
  设 \(\mu\) 和 \(\mu'\) 是两个复测度，且 \(\nu = \mu \otimes \mu'\) 。若 \(\mathbf{f}\) 是 \(\nu\) -可积的（换句话说，\(|\nu|\) -可积的），则把本定理应用于测度 \(|\mu|\) 和 \(|\mu'|\) （它们的乘积是 \(|\nu|\) ，Ch. III, §4, No.~2, 命题 3）就蕴涵：\(t' \mapsto \mathbf{f}(t, t')\) 是 \(\mu'\) -可积的（对 \(\mu\) -几乎所有的 \(t\) ）。由此，把测度 \(\mu\) 和 \(\mu'\) 分解成正测度的线性组合，即得 a) 的陈述推广到复测度。对 b) 可类似论证。
\end{enumerate}

命题 9. --- 设 F、F′ 和 G 是三个 Banach 空间，并设 \((\mathbf{x},\mathbf{y})\mapsto[\mathbf{x}\cdot\mathbf{y}]\) 是 \(F\times F'\) 到 G 中的一个连续双线性映射。设 f（相应地，\(f'\) ）是定义在 T（相应地，\(T'\) ）上、取值于 F（相应地，\(F'\) ）中且关于 \(\mu\) （相应地，\(\mu'\) ）本质可积的函数。用 g 表示函数 \((t,t')\mapsto[\mathbf{f}(t)\cdot\mathbf{f}'(t')]\) ；则 g 关于 \(\mu\otimes\mu'\) 本质可积，且

\[
\iint \left[ \mathbf {f} (t) \cdot \mathbf {f} ^ {\prime} \left(t ^ {\prime}\right) \right] d \mu (t) d \mu^ {\prime} \left(t ^ {\prime}\right) = \left[ \left(\int \mathbf {f} d \mu (t)\right) \cdot \left(\int \mathbf {f} ^ {\prime} \left(t ^ {\prime}\right) d \mu^ {\prime} \left(t ^ {\prime}\right)\right) \right].\tag{14}
\]

此外，若 \(\mathbf{f}\) 和 \(\mathbf{f}'\) 可积，则 \(\mathbf{g}\) 可积。

函数 \((t, t') \mapsto [\mathbf{f}(t) \cdot \mathbf{f}'(t')]\) 依命题 3 的推论 1 是 \((\mu \otimes \mu')\) -可测的。另一方面，若用 \(b\) 表示双线性映射 \((\mathbf{x}, \mathbf{y}) \mapsto [\mathbf{x} \cdot \mathbf{y}]\) 的范数，则

\[
\begin{array}{r l} \iint^ {\bullet} | [ \mathbf {f} (t) \cdot \mathbf {f} ^ {\prime} (t ^ {\prime}) ] | d \mu (t) d \mu^ {\prime} (t ^ {\prime}) & \leqslant b \iint^ {\bullet} | \mathbf {f} (t) | \cdot | \mathbf {f} ^ {\prime} (t ^ {\prime}) | d \mu (t) d \mu^ {\prime} (t ^ {\prime}) \\ & = b \left(\int^ {\bullet} | \mathbf {f} (t) | d \mu (t)\right) \left(\int^ {\bullet} | \mathbf {f} ^ {\prime} (t ^ {\prime}) | d \mu^ {\prime} (t ^ {\prime})\right) \end{array}
\]

这依据命题 8。这表明 \([\mathbf{f}(t) \cdot \mathbf{f}'(t')]\) 关于 \(\mu \otimes \mu'\) 本质可积（§1, No.~3, 命题 9）。设 \(\mathbf{f}\) 和 \(\mathbf{f}'\) 可积：则 \(\mathbf{f}\) 和 \(\mathbf{f}'\) 是适度的，且 \(\mathbf{g}\) 是适度的（命题 5 的推论 2），从而可积（§1, No.~3, 命题 9 的推论）。在这一情形，公式 (14) 由 Lebesgue--Fubini 定理及积分的线性性得出（Ch. IV, §4, No.~2, 定理 1）。为完成 \(\mathbf{f}\) 和 \(\mathbf{f}'\) 本质可积情形的处理，把 (14) 应用于两个可积函数 \(\mathbf{f}_1\) 和 \(\mathbf{f}_1'\) （分别局部几乎处处等于 \(\mathbf{f}\) 和 \(\mathbf{f}'\) ），并注意 \([\mathbf{f} \cdot \mathbf{f}'] = [\mathbf{f}_1 \cdot \mathbf{f}_1']\) 在 \(T \times T'\) 中局部几乎处处成立（命题 4）。

这一结果推广到复测度的乘积。

\subsection{对两测度乘积的运算}

命题 10. --- 设 \(g\) （相应地，\(g'\) ）是复函数（或取值于 \(\overline{\mathbf{R}}\) 中的函数），定义在 \(T\) （相应地，\(T'\) ）上。

\begin{enumerate}
\def\labelenumi{\alph{enumi})}
\tightlist
\item
  若 \(g\) （相应地，\(g'\) ）关于 \(\mu\) （相应地，\(\mu'\) ）局部可积，则函数 \(g \otimes g': (t, t') \mapsto g(t)g'(t')\) 关于 \(\nu = \mu \otimes \mu'\) 局部可积，且
\end{enumerate}

\[
(g \cdot \mu) \otimes (g ^ {\prime} \cdot \mu^ {\prime}) = (g \otimes g ^ {\prime}) \cdot (\mu \otimes \mu^ {\prime}).\tag{15}
\]

\begin{enumerate}
\def\labelenumi{\alph{enumi})}
\setcounter{enumi}{1}
\item
  反之，若 \(g \otimes g'\) 是局部 \(\nu\) -可积的，且 \(g'\) 不是局部 \(\mu'\) -可忽略的，则 \(g\) 是局部 \(\mu\) -可积的。
\item
  设 K 和 \(K'\) 分别是 T 和 \(T'\) 的两个紧子集；命题 8 的推论 2 表明函数 \((t, t') \mapsto g(t)g'(t')\varphi_{K \times K'}(t, t')\) （它等于 \((g\varphi_{K}) \otimes (g'\varphi_{K'})\) ）是 \(\nu\) -可积的。因此，\(g \otimes g'\) 是局部 \(\nu\) -可积的。于是立即验证 (15) 的第二端满足乘积测度的特征性质（Ch. III, §4, No.~1, 定理 1）。
\item
  现在设 \(g \otimes g'\) 是局部 \(\nu\) -可积的，且 \(g'\) 不是局部 \(\mu'\) -可忽略的。设 \(\mu_1\) 是满足 \(\mu_1 \leqslant \mu\) 的具紧支集正测度；由于 \(g \otimes g'\) 是 \((\mu_1 \otimes \mu')\) -可测的，故 \(t \mapsto g(t)g'(t')\) 是 \(\mu_1\) -可测的（除去一个局部 \(\mu'\) -可忽略的 \(t'\) 值集外）（命题 2）。
\end{enumerate}

由于 \(g'\) 局部 \(\mu'\) -几乎处处不为零，由此推出 g 是 \(\mu_{1}\) -可测的，再把 \(\mu\) 分解成具紧支集测度族之和，即得 g 是 \(\mu\) -可测的（ \(\S2\) , No.~3, 命题 4 及 \(\S2\) , No.~2, 命题 2）。证明了这一点之后，就可化归到 g 和 \(g'\) \(\geqslant0\) 的情形，如有必要则把 g 和 \(g'\) 换成它们的绝对值。设 K 是 T 的任意紧子集，并设 \(K'\) 是 \(T'\) 的满足 \(\int g'\varphi_{K'}d\mu'\neq0\) 的紧子集。由命题 8，

\[
\left(\int^ {\bullet} g \varphi_ {K} d \mu\right) \left(\int^ {\bullet} g ^ {\prime} \varphi_ {K ^ {\prime}} d \mu^ {\prime}\right) = \iint^ {\bullet} (g \otimes g ^ {\prime}) \varphi_ {K \times K ^ {\prime}} d \mu d \mu^ {\prime} <   + \infty .
\]

于是第一端的第一个因子是有限的，这就完成了证明。

由于 Ch. III, §4, No.~2 的命题 3，本命题推广到复测度。

命题 11. --- 设 \(\pi\) （相应地，\(\pi'\) ）是 T（相应地，T′）到一个局部紧空间 T \(_1\) （相应地，T \(_1'\) ）中的映射。

\begin{enumerate}
\def\labelenumi{\alph{enumi})}
\item
  若 \(\pi\) （相应地，\(\pi'\) ）是 \(\mu\) -逆紧的（相应地，\(\mu'\) -逆紧的），则映射 \(\pi \times \pi'\) 是 \((\mu \otimes \mu')\) -逆紧的，且 \((\pi \times \pi')(\mu \otimes \mu') = \pi(\mu) \otimes \pi'(\mu')\) 。
\item
  反之，若 \(\pi \times \pi'\) 是 \((\mu \otimes \mu')\) -逆紧的且 \(\mu' \neq 0\) ，则 \(\pi\) 是 \(\mu\) -逆紧的。
\item
  因为 \(\pi \times \pi'\) 依 No.~2 命题 3 的推论 1 是 \((\mu \times \mu')\) -可测的。另一方面，若 K（相应地，\(K'\) ）是 \(T_1\) （相应地，\(T_1'\) ）的紧子集，则 \(\overline{\pi}^1(K)\) 和 \(\overline{\pi}'(K')\) 分别关于 \(\mu\) 和 \(\mu'\) 本质可积，因此 \(\overline{\pi}^1(K) \times \overline{\pi}'(K')\) 关于 \(\mu \otimes \mu'\) 本质可积（命题 8 的推论 2）。这就证明了 \(\pi \times \pi'\) 是 \((\mu \times \mu')\) -逆紧的。现在置 \(\mu_1 = \pi(\mu)\) 、\(\mu_1' = \pi'(\mu')\) 、\(\nu_1 = (\pi \times \pi')(\mu \otimes \mu')\) ；对 \(f \in \mathcal{K}(T_1)\) 和 \(f' \in \mathcal{K}(T_1')\) ，有
\end{enumerate}

\[
\begin{array}{r l} \iint f (\pi (t)) f ^ {\prime} (\pi^ {\prime} (t ^ {\prime})) d \mu (t) d \mu^ {\prime} (t ^ {\prime}) \\ & = \left(\int f (\pi (t)) d \mu (t)\right) \left(\int f ^ {\prime} (\pi^ {\prime} (t ^ {\prime})) d \mu^ {\prime} (t ^ {\prime})\right) \end{array}
\]

（命题 8 的推论 2），这就证明了 \(\nu_{1} = \mu_{1}\otimes \mu_{1}^{\prime}\) （Ch. III, §4, No.~1, 定理 1）。

\begin{enumerate}
\def\labelenumi{\alph{enumi})}
\setcounter{enumi}{1}
\tightlist
\item
  现在设 \(\pi \times \pi'\) 是 \(\mu \otimes \mu'\) -逆紧的且 \(\mu' \neq 0\) 。设 \(\mu_1\) 是一个 \(\leqslant \mu\) 的具紧支集测度。由于函数 \(\pi \times \pi'\) 关于 \(\mu_1 \otimes \mu'\) 可测，故映射 \(t \mapsto (\pi(t), \pi'(t'))\) 是 \(\mu\) -可测的（除去 \(t'\) 组成局部 \(\mu'\) -可忽略集的情形外）（No.~2, 命题 2）。由于 \(\mu' \neq 0\) ，由此推出 \(\pi\) 是 \(\mu_1\) -可测的，最终得出 \(\pi\) 是 \(\mu\) -可测的（§2, No.~3, 命题 4 及 §2, No.~2, 命题 2）。剩下要证明 \(\mu^{\bullet}(f \circ \pi) < +\infty\) 对每个函数 \(f \in \mathcal{K}_{+}(\mathrm{T}_{1})\) 成立。若 \(\mu\) 为零，这一性质是明显的。若 \(\mu\) 不为零，则 \(\mu \otimes \mu'\) 也不为零，于是 \((\pi \times \pi') (\mu \otimes \mu') \neq 0\) （§6, No.~2, 命题 2 的推论 1）。由 Ch. III, §4, No.~1 的引理 1，存在两个函数 \(g \in \mathcal{K}_{+}(\mathrm{T}_{1})\) 、\(g' \in \mathcal{K}_{+}(\mathrm{T}_{1}')\) ，使得
\end{enumerate}

\[
\langle (\pi \times \pi^ {\prime}) (\mu \otimes \mu^ {\prime}), g \otimes g ^ {\prime} \rangle \neq 0.
\]

依像测度的定义，这个表达式等于 \(\langle\mu\otimes\mu',(g\circ\pi)\otimes(g'\circ\pi')\rangle\) ，故命题 8 蕴涵 \(\mu^{\bullet}(g'\circ\pi')\neq0\) 。于是由命题 8 以及 §6, No.~2 的命题 2，有

\[
\begin{array}{r l} \left(\int^ {\bullet} (f \circ \pi) d \mu\right) & \left(\int^ {\bullet} (g ^ {\prime} \circ \pi^ {\prime}) d \mu^ {\prime}\right) = \iint^ {\bullet} (f \circ \pi) \otimes (g ^ {\prime} \circ \pi^ {\prime}) d \mu d \mu^ {\prime} \\ & = \iint^ {\bullet} (f \otimes g ^ {\prime}) d ((\pi \times \pi^ {\prime}) (\mu \otimes \mu^ {\prime})) <   + \infty . \end{array}
\]

于是第一端中的第一个积分是有限的，这就完成了证明。

这一结果立即推广到两个复测度的乘积（把该陈述应用于它们的绝对值）。对下列命题同样成立。

命题 12. --- 设 X（相应地，X′）是 T（相应地，T′）的一个局部紧子空间。则诱导测度 \((\mu \otimes \mu')_{X \times X'}\) （在局部紧子空间 \(X \times X'\) ——\(T \times T'\) 的子空间——上）等于乘积 \(\mu_{X} \otimes \mu'_{X'}\) ，即 X 和 \(X'\) 上分别由 \(\mu\) 和 \(\mu'\) 诱导的测度之积。

因为若 \(f \in \mathcal{K}(\mathrm{X})\) 且 \(f' \in \mathcal{K}(\mathrm{X}')\) ，则

\[
\iint_ {\mathrm{X} \times \mathrm{X} ^ {\prime}} f (t) f ^ {\prime} (t ^ {\prime}) d \mu (t) d \mu^ {\prime} (t ^ {\prime}) = \left(\int_ {\mathrm{X}} f (t) d \mu (t)\right) \left(\int_ {\mathrm{X} ^ {\prime}} f ^ {\prime} (t ^ {\prime}) d \mu^ {\prime} (t ^ {\prime})\right)
\]

这依据命题 8 的推论 2，它依诱导测度的定义（Ch. IV, §5, No.~7）证明了

\[
(\mu \otimes \mu^ {\prime}) _ {\mathbf {X} \times \mathbf {X} ^ {\prime}} = \mu_ {\mathbf {X}} \otimes \mu_ {\mathbf {X} ^ {\prime}} ^ {\prime}
\]

（Ch. III, §4, No.~1, 定理 1）。

\subsection{关于有限个测度乘积的积分}

前面的结果可以毫无困难地推广到有限个测度的乘积。例如，设 \(T_{1}, T_{2}, T_{3}\) 是三个局部紧空间，\(\mu_{i}\) 是 \(T_{i}\) 上的正测度（i = 1, 2, 3），并设 \(\nu = \mu_{1} \otimes \mu_{2} \otimes \mu_{3}\) 是 \(T = T_{1} \times T_{2} \times T_{3}\) 上的乘积测度。设 f 是 \(\nu\) -可积、取值于 \(\overline{R}\) 或一个 Banach 空间中的函数；Lebesgue--Fubini 定理的第一次应用表明，除去点 \((t_{1}, t_{2}) \in T_{1} \times T_{2}\) （组成一个对 \(\mu_{1} \otimes \mu_{2}\) 可忽略的集）外，函数 \(t_{3} \mapsto \mathbf{f}(t_{1}, t_{2}, t_{3})\) 是 \(\mu_{3}\) -可积的，函数

\[
(t _ {1}, t _ {2}) \mapsto \int \mathbf {f} (t _ {1}, t _ {2}, t _ {3}) d \mu_ {3} (t _ {3}),
\]

（在 \(\mathrm{T}_1\times \mathrm{T}_2\) 中几乎处处有定义）是 \((\mu_{1}\otimes \mu_{2})\) -可积的，且

\[
\iiint \mathbf {f} (t _ {1}, t _ {2}, t _ {3}) d \nu (t _ {1}, t _ {2}, t _ {3}) = \iint d \mu_ {1} (t _ {1}) d \mu_ {2} (t _ {2}) \int \mathbf {f} (t _ {1}, t _ {2}, t _ {3}) d \mu_ {3} (t _ {3}).
\]

同一定理的第二次应用表明，对几乎每个 \(t_{1} \in T_{1}\) ，函数 \(t_{2} \mapsto \int \mathbf{f}(t_{1}, t_{2}, t_{3}) d\mu_{3}(t_{3})\) 在 \(T_{2}\) 中几乎处处有定义且是 \(\mu_{2}\) -可积的；此外，函数

\[
t _ {1} \mapsto \int d \mu_ {2} (t _ {2}) \int \mathbf {f} (t _ {1}, t _ {2}, t _ {3}) d \mu_ {3} (t _ {3}),
\]

（在 \(T_{1}\) 中几乎处处有定义）是 \(\mu_{1}\) -可积的，且

\[
\iiint \mathbf {f} (t _ {1}, t _ {2}, t _ {3}) d \nu (t _ {1}, t _ {2}, t _ {3}) = \int d \mu_ {1} (t _ {1}) \int d \mu_ {2} (t _ {2}) \int \mathbf {f} (t _ {1}, t _ {2}, t _ {3}) d \mu_ {3} (t _ {3}).
\]

类似地可以证明，对几乎每个 \(t_{1} \in T_{1}\) ，函数 \((t_{2}, t_{3}) \mapsto \mathbf{f}(t_{1}, t_{2}, t_{3})\) 是 \((\mu_{2} \otimes \mu_{3})\) -可积的，函数

\[
t _ {1} \mapsto \iint \mathbf {f} (t _ {1}, t _ {2}, t _ {3}) d \mu_ {2} (t _ {2}) d \mu_ {3} (t _ {3}),
\]

（几乎处处有定义）是 \(\mu_{1}\) -可积的，且

\[
\iiint \mathbf {f} (t _ {1}, t _ {2}, t _ {3}) d \nu (t _ {1}, t _ {2}, t _ {3}) = \int d \mu_ {1} (t _ {1}) \iint \mathbf {f} (t _ {1}, t _ {2}, t _ {3}) d \mu_ {2} (t _ {2}) d \mu_ {3} (t _ {3}).
\]

我们把对上面就两个测度乘积证明的其他结果作同样推广的任务留给读者。

\subsection{\texorpdfstring{应用：\(R^{n}\) 中欧氏球的测度}{应用：R to the n 中欧氏球的测度}}

设 \(\mu\) 是 \(\mathbf{R}\) 上的 Lebesgue 测度，\(\mu_n\) 是 \(\mathbf{R}^n\) 上的 Lebesgue 测度，即 \(n\) 个都等于 \(\mu\) 的因子之积。我们来计算欧氏单位球的测度 \(V_n = \mu_n(\mathbf{B}_n)\) 。由命题 7 的推论 2，

\[
\mathrm{V} _ {n} = \int_ {- 1} ^ {+ 1} \mu_ {n - 1} \big (\mathbf {B} _ {n} (z _ {n}) \big) d z _ {n}.\tag{16}
\]

现在，截口 \(\mathbf{B}_n(z_n)\) 是 \(\mathbf{R}^{n-1}\) 中由关系式 \(\sum_{i=1}^{n-1} z_i^2 \leqslant 1 - z_n^2\) 定义的子集，换句话说，它是球 \(\mathbf{B}_{n-1}\) 在比为 \(\sqrt{1 - z_n^2}\) 的位似下的像。但由命题 11 及公式

\[
\alpha \int_ {- \infty} ^ {+ \infty} f (\alpha x) d x = \int_ {- \infty} ^ {+ \infty} f (z) d z
\]

（对 \(f \in \mathcal{K}(\mathbf{R})\) ）立即得出，\(\mu_{n-1}\) 在位似 \(\mathbf{x} \mapsto \alpha \mathbf{x}\) 下的像是测度 \(\alpha^{1-n} \mu_{n-1}\) 。因此

\[
\mu_ {n - 1} \left(\mathbf {B} _ {n} \left(z _ {n}\right)\right) = \left(\sqrt {1 - z _ {n} ^ {2}}\right) ^ {n - 1} \mathrm{V} _ {n - 1}.
\]

把它代入 (16)，并作变量替换 \(z_{n} = \sin \varphi\) （其中 \(-\frac{\pi}{2} \leqslant \varphi \leqslant \frac{\pi}{2}\) ），即得

\[
\mathrm{V} _ {n} = \mathrm{V} _ {n - 1} \int_ {- \frac {\pi}{2}} ^ {+ \frac {\pi}{2}} \cos^ {n} \varphi d \varphi = 2 \mathrm{V} _ {n - 1} \int_ {0} ^ {\frac {\pi}{2}} \cos^ {n} \varphi d \varphi .\tag{17}
\]

但（FRV, Ch. VII, §1, No.~3, 公式 (20)）

\[
\int_ {0} ^ {\frac {\pi}{2}} \cos^ {m} \varphi d \varphi = \frac {1}{2} \frac {\Gamma \left(\frac {1}{2}\right) \Gamma \left(\frac {m + 1}{2}\right)}{\Gamma \left(\frac {m + 2}{2}\right)}
\]

而把上式代入关系式 (17) 并注意到 \(\Gamma(\frac{1}{2})\) 的表达式（FRV, VII, §1, No.~3, 公式 (21)），最终得到

\[
\mathrm{V} _ {n} = \frac {\pi^ {n / 2}}{\Gamma \left(\frac {n}{2} + 1\right)}.\tag{18}
\]

\chapter*{习题}
\phantomsection
\addcontentsline{toc}{chapter}{习题}
\setcounter{exercisegroup}{1}
\edef\CurrentExerciseGroup{\arabic{exercisegroup}}
\section*{\texorpdfstring{\S\CurrentExerciseGroup}{第{\CurrentExerciseGroup}节}}
\phantomsection
\addcontentsline{toc}{section}{第{\CurrentExerciseGroup}节习题}

\begin{enumerate}
\def\labelenumi{\arabic{enumi})}
\tightlist
\item
  证明：若 \(f\) 和 \(g\) 是两个数值函数 \(\geqslant 0\) ，则
\end{enumerate}

\[
\mu^ {\bullet} (\sup (f, g)) + \mu^ {\bullet} (\inf (f, g)) \leqslant \mu^ {\bullet} (f) + \mu^ {\bullet} (g)
\]

（参见 Ch. IV, §1, 习题 1）。

\begin{enumerate}
\def\labelenumi{\arabic{enumi})}
\setcounter{enumi}{1}
\tightlist
\item
  证明：对每个函数 \(f \geqslant 0\) （定义在 \(T\) 上），映射 \(\mu \mapsto \mu^{\bullet}(f)\) （\(\mathcal{M}_{+}(T)\) 到 \(\overline{\mathbf{R}}\) 中的映射）关于拟强拓扑是下半连续的（Ch. III, §1, 习题 8；参见 Ch. IV, §1, 习题 6 b)）。在什么条件下此映射连续？
\end{enumerate}

¶ 3) a) 设 \((f_{\alpha})_{\alpha \in \mathbb{A}}\) 是数值函数 \(\geqslant 0\) 组成的一个族，按关系 \(\leqslant\) 有向，且使映射 \(t \mapsto (f_{\alpha}(t))\) （\(\mathbf{T}\) 到 \(\mathbf{R}^{\mathbf{A}}\) 中的映射）是 \(\mu\) -可测的。设 \(f\) 是族 \((f_{\alpha})\) 的上包络。则 \(f\) 是 \(\mu\) -可测的，且

\[
\int^ {\bullet} f d \mu = \sup _ {\alpha \in A} \int^ {\bullet} f _ {\alpha} d \mu .\tag{1}
\]

\begin{enumerate}
\def\labelenumi{\alph{enumi})}
\setcounter{enumi}{1}
\tightlist
\item
  设 \((f_{\alpha})_{\alpha\in\mathrm{A}}\) 是定义在 T 上的数值函数 \(\geqslant0\) 组成的一个族，且具有下列性质：对 T 的每个紧子集 K 和每个 \(\varepsilon>0\) ，存在紧集 \(K'\subset K\) ，使得 \(\mu(K-K')\leqslant\varepsilon\) ，且在 \(K'\) 上所有函数 \(f_{\alpha}\) 的限制都是下半连续的。设 f 是族 \((f_{\alpha})\) 的上包络。证明 f 可测并满足关系式 (1)。
\end{enumerate}

给出一个映射 \(t \mapsto (f_{\alpha}(t))\) 的例子，它满足上述条件但不是 \(\mu\) -可测的。

\begin{enumerate}
\def\labelenumi{\arabic{enumi})}
\setcounter{enumi}{3}
\tightlist
\item
  设 \(\mathbf{T}\) 和 \(\mu\) 是 Ch. IV, §1 习题 5 中定义的局部紧空间和测度。
\end{enumerate}

\begin{enumerate}
\def\labelenumi{\alph{enumi})}
\item
  证明 \(\mu\) 是具有限支集测度序列的上确界，但 \(\mu^{*} \neq \mu^{\bullet}\) ，且 \(\mu\) 不是适度的。由此推出，当把其中的符号 \(\int^{\bullet}\) 换成 \(\int^{*}\) 时，即使集 A 可数，命题 11 也变为不真。
\item
  对每个实数 \(y\) ，用 \(f_y\) 表示化为点 \((0, y)\) （\(\mathbf{T}\) 的点）的单点集的特征函数。证明映射 \(t \mapsto (f_y(t))_{y \in \mathbf{R}}\) 是 \(\mu\) -可测的，但若 \(f = \sup_{y \in \mathbf{R}} f_y\) ，则 \(\mu^*(f) = +\infty\) 而 \(\sup_{y \in \mathbf{R}} \mu^*(f_y) = 0\) 。
\end{enumerate}

\begin{enumerate}
\def\labelenumi{\arabic{enumi})}
\setcounter{enumi}{4}
\tightlist
\item
  设 \((\mu_n)\) 是 \(T\) 上的正测度序列，拟强收敛于 \(\mu\) （在 \(\mathcal{M}(T)\) 中）（Ch. III, §1, 习题 8）。
\end{enumerate}

\begin{enumerate}
\def\labelenumi{\alph{enumi})}
\item
  证明：若映射 \(f\) （\(\mathbf{T}\) 到一个拓扑空间 \(\mathbf{G}\) 中的映射）是 \(\mu_n\) -可测的（对每个 \(n\) ），则它是 \(\mu\) -可测的（利用 Ch. IV, §1 的习题 6 b)）。
\item
  设 \(\mathbf{f}\) 是 T 到一个 Banach 空间 \(F\) 中的映射，它是本质 \(\mu_n\) -可积的（对每个 \(n\) ）。证明：若序列 \((\mu_n^\bullet (|\mathbf{f}|))\) 有界，则 \(\mathbf{f}\) 是本质 \(\mu\) -可积的（习题 2）。用一个例子说明不一定有 \(\int \mathbf{f} d\mu = \lim_{n\to \infty}\int \mathbf{f} d\mu_n\) （对 \(\mu_n\) 取紧空间上的一个点测度，使得 \(\lim \| \mu_n\| = 0\) ）；但当 \(\mathbf{f}\) 有界且具紧支集时，这一关系式成立。
\end{enumerate}

\begin{enumerate}
\def\labelenumi{\arabic{enumi})}
\setcounter{enumi}{5}
\tightlist
\item
  设 \(\mathfrak{K}\) 是 T 的紧子集组成的一个 \(\mu\) -稠密族。设 \(\mathbf{f}\) 是定义在 T 上、取值于一个 Banach 空间 F 中的函数，使得 \(\mathbf{f}\varphi_{\mathrm{K}}\) 是 \(\mu\) -可积的（对每个 \(\mathrm{K} \in \mathfrak{K}\) ）。若积分 \(\int \mathbf{f}\varphi_{\mathrm{K}} d\mu\) 沿有向集 \(\mathfrak{K}\) （按 \(\subset\) ）在 F 中有一个极限，则 \(\mathbf{f}\) 是本质 \(\mu\) -可积的。
\end{enumerate}

\setcounter{exercisegroup}{2}
\edef\CurrentExerciseGroup{\arabic{exercisegroup}}
\section*{\texorpdfstring{\S\CurrentExerciseGroup}{第{\CurrentExerciseGroup}节}}
\phantomsection
\addcontentsline{toc}{section}{第{\CurrentExerciseGroup}节习题}

\begin{enumerate}
\def\labelenumi{\arabic{enumi})}
\item
  设 \((\lambda_{\alpha})\) 是 T 上的复点测度组成的一个族。为使族 \((\lambda_{\alpha})\) 在赋有淡拓扑的向量空间 \(\mathcal{M}(\mathrm{T};\mathbf{C})\) 中可和，必须且只须族 \((|\lambda_{\alpha}|)\) 可和。
\item
  设 \((\lambda_{n})\) 是 T 上的复测度序列。为使序列 \((\lambda_{n})\) 在赋有淡拓扑的 \(\mathcal{M}(\mathrm{T};\mathbf{C})\) 中可和，必须且只须对每个函数 \(f\in \mathcal{K}(\mathrm{T})\) ，通项为 \(\lambda_{n}(f)\) 的级数绝对收敛（利用 Banach-Steinhaus 定理；参见 TVS, III, §4, No.~2, 定理 1 的推论 2）。
\item
  设 \(\mathbf{T}\) 是紧区间 \([0,1]\) （\(\mathbf{R}\) 的紧区间）；对每个整数 \(n > 0\) ，设 \(\lambda_{n}\) 是测度
\end{enumerate}

\[
f \mapsto \frac {1}{n} \int_ {0} ^ {1} f (x) \sin 2 n \pi x d x
\]

（T 上的）。证明：在赋有淡拓扑（或拟强拓扑（Ch. III, §1, 习题 8））的空间 \(\mathcal{M}(\mathrm{T})\) 中，测度序列 \((\lambda_n)\) 可和，但序列 \((|\lambda_n|)\) 不可和。（注意若 \(\alpha_n = n \cdot \lambda_n(f)\) 则 \(\sum_{n} \alpha_n^2 < +\infty\) ，这要利用 Bessel 不等式，再由 Cauchy-Schwarz 不等式由此推出 \(\sum_{n} |\lambda_n(f)| < +\infty\) 。）

\begin{enumerate}
\def\labelenumi{\arabic{enumi})}
\setcounter{enumi}{3}
\tightlist
\item
  设 \(\mu\) 是 T 上的一个正测度，\((\nu_{i})_{i\in I}\) 是 T 上的正测度组成的可和族，使得 \(\mu \leqslant \sum_{i}\nu_{i}\) 。存在 T 上的测度组成的可和族 \((\mu_{i})_{i\in I}\) ，使得 \(\mu = \sum_{i}\mu_{i}\) 且 \(\mu_{i} \leqslant \nu_{i}\) （对所有 \(i \in I\) ）。（利用命题 4，化归到 T 紧的情形，再化归到 \(I = N\) 的情形。然后用归纳法构造诸 \(\mu_{i}\) 。）
\end{enumerate}

\setcounter{exercisegroup}{3}
\edef\CurrentExerciseGroup{\arabic{exercisegroup}}
\section*{\texorpdfstring{\S\CurrentExerciseGroup}{第{\CurrentExerciseGroup}节}}
\phantomsection
\addcontentsline{toc}{section}{第{\CurrentExerciseGroup}节习题}

\begin{enumerate}
\def\labelenumi{\arabic{enumi})}
\item
  设 T 是 R 的紧区间 \([0,1]\) ，\(\mu\) 是 T 上的 Lebesgue 测度。设 X 是把区间 \([0,1]\) 赋以离散拓扑所得的局部紧空间。映射 \(t \mapsto \lambda_{t} = \varepsilon_{t}\) （T 到 \(\mathcal{M}(X)\) 中的映射）是标量本质可积的，其积分为 \(\nu = 0\) 。对 X 上等于 1 的常值函数 f，命题 3 的公式 (6) 不成立。由此推出映射 \(t \mapsto \lambda_{t}\) 不是 \(\mu\) -适宜的。
\item
  设 T 和 \(\mu\) 是 Ch. IV, §1 习题 5 中定义的局部紧空间和测度 \(\geqslant 0\) 。取 X 为集 T 赋以由 \(R^{2}\) 诱导的拓扑所得的空间，于是 X 是局部紧的且在无穷远处可数。设 \(t \mapsto \lambda_{t}\) 是 T 到 \(\mathcal{M}(X)\) 中的映射，满足 \(\lambda_{t} = \varepsilon_{t}\) （当 \(t = (0, y)\) ），且 \(\lambda_{t} = \varepsilon_{t}/n^{3}\) （当 \(t = (1/n, k/n^{2})\) ）。证明映射 \(t \mapsto \lambda_{t}\) 是 \(\mu\) -适宜的、淡 \(\mu\) -可测的、但不淡连续，且若 \(\nu = \int \lambda_{t} d\mu(t)\) 而 \(f = \varphi_{I}\) ，其中 I 是 \(R^{2}\) 中以 0 为中心、边长为 2 的正方形，则
\end{enumerate}

\[
\int^ {*} f (x) d \nu (x) <   \int^ {*} d \mu (t) \int^ {\bullet} f (x) d \lambda_ {t} (x) = + \infty .
\]

\begin{enumerate}
\def\labelenumi{\arabic{enumi})}
\setcounter{enumi}{2}
\tightlist
\item
  设 \(\mathrm{T}, \mathrm{T}'\) 是两个局部紧空间，\(\mu\) 是 \(\mathrm{T}\) 上的一个正测度，\(\mu'\) 是 \(\mathrm{T}'\) 上的一个正测度。在乘积空间 \(\mathrm{X} = \mathrm{T} \times \mathrm{T}'\) 上，置 \(\lambda_t = \varepsilon_t \otimes \mu'\) 、\(\lambda_{t'}' = \mu \otimes \varepsilon_{t'}\) （对 \(t \in \mathrm{T}\) 、\(t' \in \mathrm{T}'\) ）。映射 \(t \mapsto \lambda_t\) （相应地，\(t' \mapsto \lambda_{t'}'\) ）——\(\mathrm{T}\) （相应地，\(\mathrm{T}'\) ）到 \(\mathcal{M}(\mathrm{X})\) 中的映射——是淡连续且 \(\mu\) -适宜的（相应地，\(\mu'\) -适宜的），且
\end{enumerate}

\[
\nu = \mu \otimes \mu^ {\prime} = \int \lambda_ {t} d \mu (t) = \int \lambda_ {t ^ {\prime}} ^ {\prime} d \mu^ {\prime} (t ^ {\prime})
\]

（Ch. III, §4, No.~1；参见 §8）。取 T 为 R 的区间 {[}0,1{]}，取 \(\mu\) 为 Lebesgue 测度，取 T′ 为赋以离散拓扑的区间 {[}0,1{]}，并取 \(\mu'\) 为 T′ 上由每点质量 +1 定义的测度。设 \(f\) 是 X 的“对角线”（由点 \((t,t)\) （其中 \(t\) 取遍 {[}0,1{]}）组成之集）的特征函数。证明 \(f\) 是上半连续的，且

\[
1 = \int^ {*} d \mu (t) \int^ {*} f (x) d \lambda_ {t} (x) <   \int^ {*} f (x) d \nu (x) = + \infty
\]

以及

\[
0 = \int^ {\bullet} f (x) d \nu (x) <   \int^ {\bullet} d \mu (t) \int^ {\bullet} f (x) d \lambda_ {t} (x) = 1.
\]

¶ 4) 设空间 T、T′、X 和测度 \(\mu, \mu', \nu\) 的意义与习题 3 中相同，考虑乘积空间 Y = T × X = T × (T × T′)，以及 Y 上的乘积测度 \(\varpi = \mu \otimes \nu\) 。若对 \(t \in \mathrm{T}\) 置 \(\rho_t = \varepsilon_t \otimes \nu\) ，则有 \(\varpi = \int \rho_t d\mu(t)\) 。设 H 是 T 的一个对 \(\mu\) 不可测的子集（Ch. IV, §4, 习题 8），并设 A 是由点 \((t_1, t_2, t_3)\) ——满足条件 \(t_1 = t_3\) 、\(t_2 \in \mathrm{H}\) ——组成的 Y 的子集。证明特征函数 \(\varphi_A\) 对 \(\varpi\) 是局部可忽略的，但 \(\varphi_A\) 不是 \(\rho_t\) -可测的（对任何 \(t \in \mathrm{T}\) 的值）。

\begin{enumerate}
\def\labelenumi{\arabic{enumi})}
\setcounter{enumi}{4}
\item
  给出一个 \(\mu\) -适宜族 \(t \mapsto \lambda_t\) 和一个数值函数 \(f\) （定义在 \(X\) 上）的例子，使得 \(f\) 是 \(\lambda_t\) -可测的（对所有 \(t \in T\) ），但对测度 \(\nu = \int \lambda_t d\mu(t)\) 不可测（取 \(X = T\) 和 \(\lambda_t = \varepsilon_t\) （对所有 \(t \in T\) ））。
\item
  \begin{enumerate}
  \def\labelenumii{\alph{enumii})}
  \tightlist
  \item
    证明：在 §3 中涉及淡连续映射 \(t \mapsto \lambda_t\) 这一概念的所有结果中，都可以把这一假设换成下列假设：对每个函数 \(g \in \mathcal{K}_+(X)\) ，函数 \(t \mapsto \lambda_t(g)\) 在 \(T\) 上是下半连续的。
  \end{enumerate}
\end{enumerate}

\begin{enumerate}
\def\labelenumi{\alph{enumi})}
\setcounter{enumi}{1}
\tightlist
\item
  考虑一个标量本质可积映射 \(t \mapsto \lambda_t\) （\(T\) 到 \(\mathcal{M}_+(X)\) 中的映射），它满足下列条件：
\end{enumerate}

对 T 的每个紧子集 K 和每个 \(\varepsilon > 0\) ，存在紧集 \(K_{1} \subset K\) ，使得 \(\mu(K - K_{1}) \leqslant \varepsilon\) ，且在 \(K_{1}\) 上，所有函数 \(t \mapsto \langle f, \lambda_{t} \rangle\) （其中 f 取遍 \(\mathcal{K}_{+}(X)\) ）的限制都是下半连续的。

证明映射 \(t \mapsto \lambda_t\) 是 \(\mu\) -适宜的。

\begin{enumerate}
\def\labelenumi{\arabic{enumi})}
\setcounter{enumi}{6}
\tightlist
\item
  \begin{enumerate}
  \def\labelenumii{\alph{enumii})}
  \tightlist
  \item
    设 \(\Lambda : t \mapsto \lambda_t\) 是 T 到 \(\mathcal{M}_{+}(X)\) 中的一个标量本质可积映射，且 \(\nu = \int \lambda_t d\mu(t)\) 。证明对每个下半连续函数 \(f \geqslant 0\) （定义在 \(X\) 上），
  \end{enumerate}
\end{enumerate}

\[
\int^ {\bullet} f (x) d \nu (x) \leqslant \int^ {\bullet} d \mu (t) \int^ {\bullet} f (x) d \lambda_ {t} (x).
\]

\begin{enumerate}
\def\labelenumi{\alph{enumi})}
\setcounter{enumi}{1}
\tightlist
\item
  设 \(\Lambda\) 是 \(\mu\) -预适宜的，用 \(\mu'\) 表示一个正测度 \(\leqslant \mu\) ，并写 \(\mu = \mu' + \mu''\) 、\(\nu' = \int \lambda_t d\mu'(t)\) 。由 \(a\) ) 推出，对每个下半连续函数 \(f \geqslant 0\) （\(\nu\) -可积的），
\end{enumerate}

\[
\int^ {\bullet} f (x) d \nu^ {\prime} (x) = \int^ {\bullet} d \mu^ {\prime} (t) \int^ {\bullet} f (x) d \lambda_ {t} (x).
\]

把这一结果推广到下半连续函数 \(f \geqslant 0\) （\(\nu\) -适度的）。由此推出，若 \(\nu\) 是一个适度测度（特别地，若 X 在无穷远处可数），则 \(\Lambda\) 是 \(\mu\) -适宜的。

\begin{enumerate}
\def\labelenumi{\arabic{enumi})}
\setcounter{enumi}{7}
\tightlist
\item
  设 \(\Lambda : t \mapsto \lambda_t\) 是 \(T\) 到 \(\mathcal{M}_+(X)\) 中的一个映射。
\end{enumerate}

\begin{enumerate}
\def\labelenumi{\alph{enumi})}
\item
  设 \(\mu_{1},\ldots,\mu_{n}\) 是 T 上的测度，且 \(\mu=\mu_{1}+\cdots+\mu_{n}\) 。为使 \(\Lambda\) 是 \(\mu\) -适宜的，必须且只须 \(\Lambda\) 对每个 i 都是 \(\mu_{i}\) -适宜的（利用分解引理）。
\item
  设 \(\mu\) 是递增有向族 \((\mu_i)_{i \in I}\) 的上确界。为使 \(\Lambda\) 是 \(\mu\) -适宜的，必须且只须 \(\Lambda\) 是 \(\mu_i\) -适宜的（对所有 \(i \in I\) ）且是标量本质 \(\mu\) -可积的。
\item
  设 \(\mu\) 是正测度组成的可和族 \((\mu_j)_{j \in \mathbb{J}}\) 之和。为使 \(\Lambda\) 是 \(\mu\) -适宜的，必须且只须 \(\Lambda\) 是 \(\mu_j\) -适宜的（对所有 \(j \in \mathbb{J}\) ），且 \(\Lambda\) 是标量本质 \(\mu\) -可积的。
\item
  为使 \(\Lambda\) 是 \(\mu\) -适宜的，必须且只须 \(\Lambda\) 是标量本质 \(\mu\) -可积的，且 \(\Lambda\) 是 \(\mu'\) -预适宜的（对每个满足 \(\mu' \leqslant \mu\) 的具紧支集测度）。
\end{enumerate}

\begin{enumerate}
\def\labelenumi{\arabic{enumi})}
\setcounter{enumi}{8}
\tightlist
\item
  \begin{enumerate}
  \def\labelenumii{\alph{enumii})}
  \tightlist
  \item
    设 \(\mathrm{T}\) 和 \(\mathrm{X}\) 是两个局部紧空间，\(\mathcal{C}_{+}(\mathrm{T})\) 是 \(\mathrm{T}\) 上所有正连续函数组成的凸锥，而 \(\mathrm{V}\) 是 \(\mathcal{H}_{+}(\mathrm{T})\) 到 \(\mathcal{C}_{+}(\mathrm{T})\) 中的一个映射，具有下列性质：
  \end{enumerate}
\end{enumerate}

\[
\begin{array}{r l} \mathrm{V} (f + g) = \mathrm{V} f + \mathrm{V} g & \text {if} f, g \in \mathcal {K} _ {+} (\mathrm{X}) \\ \mathrm{V} (t f) = t (\mathrm{V} f) & \text {if} f \in \mathcal {K} _ {+} (\mathrm{X}), t \in \mathbf {R} _ {+}. \end{array}
\]

证明存在唯一的扩散 \(\Lambda\) （\(T\) 到 \(X\) 中的扩散），使得 \(\Lambda f = Vf\) （对所有 \(f \in \mathcal{K}_{+}(X)\) ）。证明：若把锥 \(\mathcal{C}_{+}(T)\) 换成 \(T\) 上正的且局部有界的下半连续函数组成的锥，这一结果仍成立（参见习题 6）。

\begin{enumerate}
\def\labelenumi{\alph{enumi})}
\setcounter{enumi}{1}
\tightlist
\item
  设 X 的拓扑有一个可数基。证明：若把锥 \(\mathcal{C}_{+}(\mathrm{T})\) 换成 T 上正的且局部有界的普遍可测函数组成的锥，上述结果仍成立。
\end{enumerate}

\begin{enumerate}
\def\labelenumi{\arabic{enumi})}
\setcounter{enumi}{9}
\item
  设 \(\mathrm{T}, \mathrm{X}, \mathrm{Y}\) 是三个在无穷远处可数的局部紧空间，\(\Lambda : t \mapsto \lambda_t\) 是 \(\mathrm{T}\) 到 \(\mathrm{X}\) 中的一个扩散，\(\mathrm{H}: x \mapsto \eta_x\) 是 \(\mathrm{X}\) 到 \(\mathrm{Y}\) 中的一个扩散；称 \(\Lambda\) 与 \(\mathrm{H}\) 是可复合的，如果函数 \(\Lambda(\mathrm{Hg})\) 在 \(\mathrm{T}\) 上局部有界（对每个 \(g \in \mathcal{K}_+(\mathrm{Y})\) ）。证明：若 \(\Lambda\) 与 \(\mathrm{H}\) 可复合，则 \(\lambda_t\) 属于 \(\mathrm{H}\) 的定义域（对所有 \(t \in \mathrm{T}\) ），且映射 \(t \mapsto \lambda_t\mathrm{H}\) 是 \(\mathrm{T}\) 到 \(\mathrm{Y}\) 中的一个扩散。把公式 (15)（命题 13）推广到这一情形。
\item
  设 \(t \mapsto \lambda_t\) 是一个 \(\mu\) -适宜映射（\(T\) 到 \(\mathcal{M}_+(X)\) 中的映射），满足下列条件：对每个紧子集 \(K\) （\(T\) 的子集），存在紧子集 \(L_K\) （\(X\) 的子集），使得 \(\operatorname{Supp}(\lambda_t) \subset L_K\) （对所有 \(t \in K\) ）。
\end{enumerate}

\begin{enumerate}
\def\labelenumi{\alph{enumi})}
\tightlist
\item
  证明：对每个函数 \(f \geqslant 0\) （定义在 \(X\) 上），
\end{enumerate}

\[
\int^ {\bullet} f (x) d \nu (x) \geqslant \int^ {\bullet} d \mu (t) \int^ {*} f (x) d \lambda_ {t} (x)
\]

（利用命题 3 a)）。

\begin{enumerate}
\def\labelenumi{\alph{enumi})}
\setcounter{enumi}{1}
\item
  若 \(f \geqslant 0\) 是局部 \(\nu\) -可忽略的，则由这样的 \(t \in \mathbb{T}\) 组成之集——使 \(f\) 不是局部 \(\lambda_t\) -可忽略的——是局部 \(\mu\) -可忽略的。
\item
  若 \(f\) 是一个 \(\nu\) -可测映射（\(X\) 到一个拓扑空间 \(X\) 中的映射），则由这样的 \(t \in T\) 组成之集——使 \(f\) 不是 \(\lambda_t\) -可测的——是局部 \(\mu\) -可忽略的。
\item
  对每个 \(\nu\) -可测函数 \(f \geqslant 0\) ，映射 \(t \mapsto \int^{*} f d\lambda_{t}\) 是 \(\mu\) -可测的，且
\end{enumerate}

\[
\int^ {\bullet} f (x) d \nu (x) = \int^ {\bullet} d \mu (t) \int^ {*} f (x) d \lambda_ {t} (x).
\]

\begin{enumerate}
\def\labelenumi{\alph{enumi})}
\setcounter{enumi}{4}
\tightlist
\item
  设 \(\mathbf{f}\) 是定义在 \(X\) 上、取值于一个 Banach 空间或 \(\overline{\mathbf{R}}\) 中的本质 \(\nu\) -可积函数。证明：由这样的 \(t \in T\) 组成之集——使 \(\mathbf{f}\) 不是 \(\lambda_t\) -可积的——是局部 \(\mu\) -可忽略的，函数 \(t \mapsto \int \mathbf{f}(x) d\lambda_t(x)\) （对 \(\mu\) 局部几乎处处有定义）是本质 \(\mu\) -可积的，且
\end{enumerate}

\[
\int \mathbf {f} (x) d \nu (x) = \int d \mu (t) \int \mathbf {f} (x) d \lambda_ {t} (x).
\]

\setcounter{exercisegroup}{4}
\edef\CurrentExerciseGroup{\arabic{exercisegroup}}
\section*{\texorpdfstring{\S\CurrentExerciseGroup}{第{\CurrentExerciseGroup}节}}
\phantomsection
\addcontentsline{toc}{section}{第{\CurrentExerciseGroup}节习题}

\begin{enumerate}
\def\labelenumi{\arabic{enumi})}
\item
  设 T 和 X 是两个紧空间，\(\pi\) 是 T 到 X 中的一个连续映射，g 是定义在 T 上的连续且有限的数值函数。对每个 \(x \in X\) ，用 \(f(x)\) 表示 \(g(t)\) 在集 \(\overline{\pi}^{-1}(x)\) 上的下确界。给出一个 f 不连续的例子。（取 T = R 中的 {[}0,1{]}，对 \(\pi\) 取把 T 中 0 与 1 等同起来所得商空间上 T 的典范映射。）
\item
  设 X 和 \(\nu\) 是 Ch. IV, §1 习题 5 中定义的局部紧空间和测度；设 f 是 X 中由点 \((0,y)\) 组成的集 D 的特征函数。证明：在下列两种情形下，\(\nu=\int g(t)\varepsilon_{\pi(t)}d\mu(t)\) ，且
\end{enumerate}

\[
\int^ {*} f (\pi (t)) g (t) d \mu (t) <   \int^ {*} f (x) d \nu (x) = + \infty .
\]

\begin{enumerate}
\def\labelenumi{\alph{enumi})}
\tightlist
\item
  取 T = X，对 \(\mu\) 取由质量 \(\log n/n^{3}\) （在每点 \((1/n, k/n^{2})\) 处）所定义的测度。对 \(\pi\) 取恒等映射，g 由下列条件定义：
\end{enumerate}

\[
g (0, y) = 0, \quad g (1 / n, k / n ^ {2}) = 1 / \log n;
\]

\(g\) 是连续的，但 \(g(t) > 0\) 并非对每个 \(t \in \mathbf{T}\) 成立。

\begin{enumerate}
\def\labelenumi{\alph{enumi})}
\setcounter{enumi}{1}
\tightlist
\item
  取 T 为赋以离散拓扑的集 X，对 \(\mu\) 取由与 \(\nu\) 相同的质量定义的测度。取 \(g(t) = 1\) （对所有 \(t\) ），对 \(\pi\) 取恒等映射：后者是连续的，但不是逆紧的。
\end{enumerate}
